\documentclass[10pt,a4paper]{amsart}
\usepackage[margin=19mm]{geometry}
\usepackage{amsmath,amssymb,mathtools}
\usepackage[scr=rsfs,cal=euler]{mathalfa}
\usepackage{xfrac}
\usepackage[unicode,hidelinks,bookmarks=false]{hyperref}
\newcommand{\ie}{i.e.\ }
\newcommand{\cf}{cf.\ }
\newcommand{\resp}{respectively\ }
\newcommand{\Subsection}{\subsection}
\providecommand{\nobreakdash}{\nobreak\hbox{-}}
\setlength{\emergencystretch}{2em}
\multlinegap0pt
\usepackage{bm}
\usepackage{dsfont}
\usepackage{xcolor}
\usepackage{amssymb}
\newtheorem{lemma}{Lemma}
\newcommand{\lV}{\lVert}
\newcommand{\lf}{\left}
\newcommand{\rV}{\rVert}
\newcommand{\rg}{\right}
\title{Existence and uniqueness of solutions of unsteady Darcy-Brinkman problem for modelling miscible reactive flows in porous media}
\author{Pankaj Roy, Satyajit Pramanik}
\date{Source: arXiv:2506.22225v2 (CC BY 4.0)}
\begin{document}
\maketitle

\noindent\emph{Source and licence.} Existence and uniqueness of solutions of unsteady Darcy-Brinkman problem for modelling miscible reactive flows in porous media. Version arXiv:2506.22225v2 (CC BY 4.0), distributed by arXiv under the Creative Commons Attribution 4.0 International licence, http://creativecommons.org/licenses/by/4.0/, as recorded in the arXiv licence metadata for this exact version and shown on the abstract page https://arxiv.org/abs/2506.22225v2. Full licence notice: \texttt{licenses/arxiv-2506.22225-CC-BY-4.0.txt}.\par\smallskip

\noindent\textit{Editorial scope.} Counted unit: the article's lemma lemma:L4 (source0.tex lines 461--463). The article's complete original proof of it is delivered with it (source0.tex lines 465--500, one \texttt{proof} environment of 36 lines). No mathematics is written, repaired or re-derived: the statement and the proof are the article's own lines, in the article's own notation, and no retained line is substituted or re-flowed.  Retained uncounted as the source-local context the proof needs: lines 308--314; lines 380--382.  Nothing was omitted from any retained span, so the package carries no omission record.  No retained line contains a citation, so the wrapper carries no bibliography and no bibliography entry.  No retained line contains a figure environment, an image-inclusion command or drawing code, so no image is delivered and the package declares no asset dependency.  Editorial formatting only, in addition: an AMS \texttt{amsart} preamble that restates the article's own declarations and macros used by the retained lines, namely the environments \texttt{lemma} on separate counters, together with the standard packages those lines need.  The retained environments print in this document as Lemma 1 in order of appearance, whereas the article prints the same items with the numbering of its own full document.  The complete native source of the article as distributed in the arXiv e-print for this exact version is bundled unchanged as \texttt{sources/180612/source0.tex}.
\begin{align}
	\left(\frac{\partial C_n\left(t\right)}{\partial t},z_j\right) & + d\left(\boldsymbol{\nabla} C_n\left(t\right), \boldsymbol{\nabla} z_j\right)+\left(\mathbf{u}_n\left(t\right) \cdot \boldsymbol{\nabla} C_n\left(t\right),z_j\right)\nonumber \\ 
    & +\kappa \left(C_n\left(t\right)\left(1-C_n\left(t\right)\right),z_j\right) = 0 \quad \forall z_j \in \left(\mathscr V_2\right)_n \quad a.e. \ on \ \left(0,T\right), \label{eq: fdC} \\
	\left \langle  \frac{\partial \mathbf{u}_n\left(t\right)}{\partial t},\mathbf w_j \right \rangle & + \mu_e\left(\boldsymbol{\nabla} \mathbf{u}_n\left(t\right), \boldsymbol{\nabla} \mathbf w_j\right) + \left( F\left(C_n\left(t\right)\right)
 \mathbf{u}_n\left(t\right),\mathbf w_j\right) \nonumber \\ 
 & -\left \langle \boldsymbol{\nabla} \cdot \mathbf{T}\left(C_n\left(t\right)\right), \mathbf w_j \right \rangle - \left(\mathbf f\left(t\right),\mathbf w_j\right) = 0 \quad \forall \mathbf w_j \in \left(\mathscr V_1\right)_n \quad a.e. \ on \ \left(0,T\right)  \label{eq: fdu},
 \end{align} 

\begin{lemma}
	$\mathbf{u}_n$ is uniformly bounded in $L^\infty\left(0,T; \mathscr S_1\right)\cap L^2\left(0,T;\mathscr V_1\right)$ and $C_n$ is uniformly bounded in $L^\infty\left(0,T; \mathscr S_2\right)\cap L^2\left(0,T;\mathscr V_2\right)$.\label{lemma:L2}
\end{lemma}

\begin{lemma}
	$ \displaystyle \frac{\partial \mathbf{u}_n}{\partial t}$ is uniformly bounded in $L^{4/N}\left(0,T;\mathscr V_1^*\right)$ for $N=2,3$. \label{lemma:L4}
\end{lemma}

\begin{proof}
		Let $\mathbf w\in \mathscr V_1$ be any non-zero vector. 
  Then $\mathbf w$ can be written as
  \begin{align*}
      \mathbf w=\mathbf w^1+\mathbf w^2\quad \text{where}\ \mathbf w^1\in \left(\mathscr V_1\right)_n,\ \mathbf w^2\in \left(\mathscr V_1\right)_n ^ \perp.
 \end{align*}
 It can be shown that
 \begin{align*}
     \left \langle \frac{\partial \mathbf{u}_n\left(t\right)}{\partial t},\mathbf w^2 \right \rangle = 0.
 \end{align*}
 Then from equation \eqref{eq: fdu} we get,
 \begin{align}
     \left \lvert \left \langle \frac{\partial \mathbf{u}_n\left(t\right)}{\partial t},\mathbf w \right \rangle \right \rvert 
     & = \lf| \left \langle \frac{\partial \mathbf{u}_n\left(t\right)}{\partial t},\mathbf w^1 \right \rangle \rg| \nonumber \\ 
     & = \left| \mu_e\left(\boldsymbol{\nabla} \mathbf{u}_n\left(t\right), \boldsymbol{\nabla} \mathbf w^1\right) +  \left(F\left(C_n\left(t\right)\right) \mathbf{u}_n\left(t\right),\mathbf w^1\right)-\left \langle \boldsymbol{\nabla} \cdot \mathbf{T}\left(C_n\left(t\right)\right) , \mathbf w^1 \right \rangle  + \left(\mathbf f\left(t\right),\mathbf w^1\right) \right| \quad  a.e. \ on \ \left(0,T\right) \nonumber\\ 
     & \le \left| \mu_e\left(\boldsymbol{\nabla} \mathbf{u}_n\left(t\right), \boldsymbol{\nabla} \mathbf w^1\right)\right| + \left|\left(  F\left(C_n\left(t\right)\right) \mathbf{u}_n\left(t\right),\mathbf w^1\right)\right|+\left|\left \langle \boldsymbol{\nabla} \cdot \mathbf{T}\left(C_n\left(t\right)\right) , \mathbf w^1 \right \rangle\right| \nonumber\\ 
     & \qquad\qquad\qquad\qquad + \left|\left(\mathbf f\left(t\right),\mathbf w^1\right)\right| \quad  a.e. \ on \ \left(0,T\right). 
 \end{align} 
 The term $\left \lvert \left(  F\left(C_n\left(t\right)\right) \mathbf{u}_n\left(t\right),\mathbf w^1\right) \right \rvert$ can be estimated as,
\begin{align}
     \left\lvert \left(  F\left(C_n\left(t\right)\right) \mathbf{u}_n\left(t\right),\mathbf w^1\right) \right\rvert &= \left \lvert \left( \mathbf{u}_n\left(t\right), \left(  F\left(C_n\left(t\right)\right) \mathbf w^1\right)\right) \right \rvert \nonumber\\& \le \|\mathbf u_n\left(t\right)\|_{L^2}\|{F\left(C_n\left(t\right)\right)}\|_{L^4}\|\mathbf w^1\|_{L^4} \nonumber\\
     &\leq M \lVert {F \left( C_n \left( t \right) \right)} \rVert_{H^1} \lVert \mathbf w^1 \rVert_{H^1} \quad a.e. \ on \ \left( 0, T \right). \label{eq: estimate of F(C_n)} 
\end{align}
Now applying Cauchy-Schwarz inequality in the rest of the terms of the right-hand side and taking $supremum$ over all $\mathbf w \in \mathscr V_1$ with $\|\mathbf w\|_{H^1} \le1$, we get,
\begin{align}
     \left \lVert \frac{\partial \mathbf{u}_n\left(t\right)}{\partial t} \right\rVert_{\mathscr V_1^*} \le \mu_e\|\boldsymbol{\nabla} \mathbf{u}_n\left(t\right)\|_{L^2} 
     & + M\|{ F\left(C_n\left(t\right)\right)}\|_{H^1} + \| \boldsymbol{\nabla} \cdot \mathbf{T}\left(C_n\left(t\right)\right) \|_{\mathscr V_1^*} \nonumber \\
     & + \| \mathbf f\left(t\right)\|_{L^2}  \quad  a.e. \ on \ \left(0,T\right)\label{eq: L51}. 
\end{align}
Here, $\nabla \cdot \mathbf{T}(C)$ is a sum of expressions of the form $\lambda D_i(D_j C \, D_k C)$. We then have the following estimate
\begin{align}\label{est: korteweg stress}
\| D_i(D_j C \, D_k C) \|_{\mathscr  V_1^*} & \leq \| D_j C \, D_k C \|_{L^2} \leq \| D_j C \|_{L^4} \, \| D_k C \|_{L^4}  \nonumber\\
\leq & M  \lV D_j C\rV^{(4-N)/4}_{L^2}  \lV D_k C\rV^{(4-N)/4}_{L^2} \| D_j C \|_{H^1}^{N/4} \, \| D_k C \|_{H^1}^{N/4} \quad\text{for}\ N = 2,3 .
\end{align}
The remaining terms on the right-hand side are uniformly bounded as a consequence of Lemma \ref{lemma:L2}.  Hence, $ \displaystyle \frac{\partial \mathbf u_n}{\partial t}$ is uniformly bounded in $L^{4/N}\left(0,T;\mathscr V_1^*\right)$ for $N=2,3$.
\end{proof}

\end{document}
